\section{Beyond Gaussian influence: higher cumulants and spin-bath corrections}
\label{sec:beyond_gaussian}

The preceding sections derived the Hamiltonian of mean force by
exploiting two simplifications simultaneously: the bath was Gaussian
(harmonic oscillators with linear coupling), and the system operators
$H_Q$ and $f$ commuted. Together these ensured that the
Hubbard--Stratonovich auxiliary field decoupled the influence functional
and that the resulting time-ordered exponential collapsed to an ordinary
exponential.

Two logically independent obstructions can prevent such a collapse in
general: (i)~noncommutativity of $H_Q$ and $f$, which forces
nontrivial time ordering, and (ii)~non-Gaussian bath statistics, which
generate higher-order connected correlators beyond the two-point kernel
$K(\tau-\tau')$. The previous section addressed obstruction~(i). Here
we address obstruction~(ii), showing how it enters the formalism through
an ordered-cumulant expansion and identifying when it can be neglected.

\subsection{Ordered-cumulant expansion}

Consider the bath trace that generates the influence functional. For a
general bath---not necessarily harmonic---the coupling operator
$B:=\sum_k c_k x_k$ from \eqref{eq:HI_model} need not produce Gaussian
statistics. The averaged time-ordered exponential admits a formal
cumulant expansion~\cite{Kubo1962,Kubo1963}
\begin{widetext}
\begin{equation}
    \left\langle
        \hat{\mathcal T}\exp\!\left[\int_0^\beta d\tau\,B(\tau)\,f\right]
    \right\rangle_{\!X}
    =
    \exp\!\left\{
        \sum_{n=1}^{\infty}\frac{1}{n!}
        \int_0^\beta \!d\tau_1\cdots\!\int_0^\beta \!d\tau_n\;
        \langle\!\langle
            \hat{\mathcal T}\,B(\tau_1)\cdots B(\tau_n)
        \rangle\!\rangle_{\!X}\;
        f^n
    \right\},
    \label{eq:ordered_cumulant}
\end{equation}
\end{widetext}
where $\langle\!\langle\cdots\rangle\!\rangle_X$ denotes the
\emph{connected} (cumulant) part of the bath correlation function,
$\langle\cdot\rangle_X=\Tr_X[\,\cdot\,e^{-\beta H_X}]/Z_X$, and
$B(\tau)=e^{\tau H_X}B\,e^{-\tau H_X}$ is the interaction-picture bath
operator. In contrast to the scalar cumulant theorem, the
time-ordering operator $\hat{\mathcal T}$ must in general be retained
inside each cumulant; the resulting objects are sometimes called
\emph{ordered cumulants}~\cite{Fox1976,Fox1979,Kubo1962}. However,
when $B(\tau)$ is a $c$-number process (as after the
Hubbard--Stratonovich transformation of
Sec.~\ref{sec:influence_to_hmf}), the ordered cumulants reduce to
ordinary scalar cumulants and the expansion becomes exact without
ordering subtleties~\cite{Tanimura2006}.

\subsection{Gaussian truncation}

For a harmonic bath with linear coupling, the collective variable
$B(\tau)$ is a Gaussian process: all connected correlators of order
$n\geq 3$ vanish identically. The cumulant
series~\eqref{eq:ordered_cumulant} therefore terminates at $n=2$:
\begin{equation}
    \left\langle
        \hat{\mathcal T}\,e^{\int_0^\beta d\tau\,B(\tau)\,f}
    \right\rangle_{\!X}
    =
    \exp\!\left[
        \frac{1}{2}\int_0^\beta\!d\tau\!\int_0^\beta\!d\tau'\,
        K(\tau-\tau')\,f^2
    \right],
    \label{eq:gaussian_truncation}
\end{equation}
recovering exactly the kernel $K(\tau-\tau')$ of
\eqref{eq:kernel_spectral} and the variance $C(\beta)$ of
\eqref{eq:Cbeta_def}. The entire influence-functional construction of
Sec.~\ref{sec:influence_to_hmf}---including the closed-form
mean-force Hamiltonian \eqref{eq:HMF_commuting_final}---is therefore
the $n\leq 2$ truncation of the general cumulant expansion.

\subsection{Leading non-Gaussian correction}

When the bath is not Gaussian, the first correction arises from the
fourth-order connected correlator (the third order vanishes for
any bath with time-reversal symmetry or vanishing odd moments).
Denoting the connected four-point function
\begin{equation}
    K_4(\tau_1,\tau_2,\tau_3,\tau_4)
    :=
    \langle\!\langle
        B(\tau_1)\,B(\tau_2)\,B(\tau_3)\,B(\tau_4)
    \rangle\!\rangle_{\!X},
    \label{eq:K4_def}
\end{equation}
the leading non-Gaussian correction to the influence functional takes
the form
\begin{widetext}
\begin{equation}
    \Phi^{(4)}
    =
    \frac{1}{4!}
    \int_0^\beta\!d\tau_1\!\int_0^\beta\!d\tau_2\!
    \int_0^\beta\!d\tau_3\!\int_0^\beta\!d\tau_4\;
    K_4(\tau_1,\tau_2,\tau_3,\tau_4)\;
    f\!\left(q(\tau_1)\right)f\!\left(q(\tau_2)\right)
    f\!\left(q(\tau_3)\right)f\!\left(q(\tau_4)\right).
    \label{eq:Phi4}
\end{equation}
\end{widetext}
This enters the exponent of the influence functional additively:
\begin{equation}
    S_{\mathrm{eff}}[q]
    =
    S_Q[q]
    -\underbrace{\tfrac{1}{2}\!\iint K\,f\,f}_{\text{Gaussian}}
    -\underbrace{\Phi^{(4)}}_{\text{leading correction}}
    -\;\cdots
    \label{eq:Seff_cumulant}
\end{equation}
Each term in the cumulant series has a transparent diagrammatic
interpretation: the $n$-th order contribution corresponds to the
sum of all connected diagrams with $n$ external legs attached to
the system coordinate at imaginary times $\tau_1,\ldots,\tau_n$.
For Gaussian baths the only connected diagram is the single
propagator line (the two-point function), and all higher connected
diagrams vanish. Non-Gaussian statistics generate connected $n$-point
vertices that are absent in the harmonic theory~\cite{FunoIshizaki2024}.

In the commuting sector, $[H_Q,f]=0$, and $f$ may be treated as a
$c$-number in the simultaneous eigenbasis of $H_Q$ and $f$. The
non-Gaussian correction then modifies the mean-force Hamiltonian in a
particularly transparent way. Defining the integrated fourth cumulant
\begin{equation}
    \kappa_4(\beta)
    :=
    \frac{1}{4!}\int_0^\beta\!d\tau_1\cdots\!\int_0^\beta\!d\tau_4\;
    K_4(\tau_1,\ldots,\tau_4),
    \label{eq:kappa4_def}
\end{equation}
the corrected mean-force Hamiltonian becomes
\begin{equation}
    H_{\mathrm{MF}}(\beta)
    =
    H_Q
    - \lambda\,f^2
    - \frac{\kappa_4(\beta)}{\beta}\,f^4
    + \cdots
    + \mathrm{const}(\beta)\,\mathbf 1.
    \label{eq:HMF_non_gaussian}
\end{equation}
Crucially, while the Gaussian correction $-\lambda f^2$ is temperature
independent (since $C(\beta)=2\beta\lambda$), the non-Gaussian
coefficient $\kappa_4(\beta)/\beta$ generically depends on temperature.
Non-Gaussian bath statistics can therefore introduce genuine
temperature dependence into the \emph{operator} content of the
mean-force Hamiltonian, even in the commuting sector where time ordering
is trivial. This marks a qualitative distinction from the purely
Gaussian case.


\subsection{Spin baths and the Gaussianization limit}
\label{sec:spin_baths}

A particularly important class of non-Gaussian environments consists
of baths of localised spins, which arise naturally in solid-state
settings (nuclear spins, paramagnetic impurities, two-level
fluctuators)~\cite{ProkofevStamp2000}. Whether a spin bath
behaves effectively as a Gaussian environment depends on how the
couplings scale with the number of bath constituents.

Consider first $N$ spin-$\tfrac{1}{2}$ degrees of freedom with
couplings $c_k\sim 1/\sqrt{N}$, chosen to discretise a smooth spectral
density $J(\omega)$. In this regime the collective bath variable
$B=\sum_k c_k\sigma_k^z$ satisfies a central limit theorem as
$N\to\infty$: its connected correlators of order $n\geq 3$ scale as
$N^{1-n/2}\to 0$ and only the two-point function
survives~\cite{SuarezSilbey1991,Makri1999}. The resulting effective
Gaussian bath is characterised by the same influence kernel
$K(\tau-\tau')$, and the mean-force construction of
Sec.~\ref{sec:influence_to_hmf} applies without modification.

When the couplings instead remain $\mathcal O(1)$ as $N$
grows---the generic situation for a physical spin bath---the central
limit theorem does not apply. Higher connected correlators scale as
$\mathcal O(1)$ and contribute at leading order. In particular the
four-point cumulant $K_4$ (and higher) do not vanish, and the Gaussian
truncation introduces an uncontrolled
error~\cite{ProkofevStamp2000,HsiehCao2018I,HsiehCao2018II}.
Even in this regime, however, an effective Gaussian description can
be recovered within linear response---that is, when the bath is only
weakly perturbed from equilibrium. There the two-point correlator
determines the full response via a fluctuation--dissipation relation,
and higher cumulants enter only at higher orders in the system--bath
interaction
strength~\cite{SuarezSilbey1991,Makri1999,YingEtAl2024}.
Beyond linear response, the non-Gaussian character of the spin bath
reasserts itself and must be treated
explicitly~\cite{HsiehCao2018II,HalataeiEtAl2025}.

The cumulant expansion~\eqref{eq:ordered_cumulant} therefore provides
a unified framework in which these distinctions become systematic. For
a given bath, the programme is straightforward: compute or bound the
connected correlators
$\langle\!\langle B(\tau_1)\cdots B(\tau_n)\rangle\!\rangle_X$. If
all $n\geq 3$ cumulants are negligible, the Gaussian construction of
Sec.~\ref{sec:influence_to_hmf} applies directly and the mean-force
Hamiltonian is given by \eqref{eq:HMF_commuting_final}. Otherwise,
$\Phi^{(4)}$ and, if needed, higher terms enter as perturbative
corrections to the effective action~\eqref{eq:Seff_cumulant}. In the
commuting sector these produce the additional operator corrections
$f^4, f^6,\ldots$ appearing in
\eqref{eq:HMF_non_gaussian}, each with a coefficient that is
determined entirely by the connected bath correlator of the
corresponding order.

\subsection{Nested Hubbard--Stratonovich representation}
\label{sec:nested_HS}

The cumulant expansion identifies the corrections; the quenched
framework of Sec.~\ref{sec:quenched} can absorb them. The essential
observation is that each higher-order term in the effective
action~\eqref{eq:Seff_cumulant} is multilocal in imaginary time and can
be reduced to a local-in-$\tau$ operator picture by a nested
Hubbard--Stratonovich transformation, at the price of promoting the
covariance of the auxiliary field to a stochastic object.

Consider the quartic correction $\Phi^{(4)}$. It is quadratic in the
composite field $\Psi(\tau,\tau'):=f(q(\tau))\,f(q(\tau'))$. Introduce a
zero-mean Gaussian random kernel $\delta K(\tau,\tau')$ with covariance
\begin{equation}
    \bigl\langle
        \delta K(\tau_1,\tau_2)\,\delta K(\tau_3,\tau_4)
    \bigr\rangle
    =
    \tfrac{1}{3}\,K_4(\tau_1,\tau_2,\tau_3,\tau_4).
    \label{eq:deltaK_covariance}
\end{equation}
The construction requires $K+\delta K$ to be positive semi-definite for the conditional Gaussian measure to be well defined. This holds perturbatively whenever the non-Gaussian correction is small, $\delta K\ll K$. For sub-Gaussian baths---including bounded spin baths where $K_4$ is typically negative---the identity survives via analytic continuation (where the auxiliary field $\delta K$ becomes complex), though the direct stochastic sampling interpretation is lost.
Since the functional $Y:=\frac{1}{2}\!\iint \delta K\,f\,f$ is linear
in the Gaussian field $\delta K$, it is itself Gaussian, with variance
$\frac{1}{8}\!\iiiint\!\frac{1}{3}K_4\,f^4 = \Phi^{(4)}$.
Consequently
\begin{align}
    e^{\Phi^{(4)}}
    &=
    \mathbb E_{\delta K}\!\left[
        \exp\!\left(
            \tfrac{1}{2}\!\iint\!\delta K(\tau,\tau')\,
            ff'\,d\tau\,d\tau'
        \right)
    \right],
    \label{eq:quartic_HS}
\end{align}
Combining this with the Gaussian kernel gives
\begin{align}
    e^{\frac{1}{2}\iint K\,ff\,+\,\Phi^{(4)}}
    &=
    \mathbb E_{\delta K}\!\left[
        e^{\frac{1}{2}\iint(K+\delta K)\,ff}
    \right].
    \label{eq:combined_kernel}
\end{align}
The exponent on the right is again bilocal and can be linearised by a
second, standard Hubbard--Stratonovich transformation: for each
realisation of $\delta K$, introduce $\xi(\tau)$ drawn from a Gaussian
with the stochastic covariance
\begin{equation}
    \xi\,|\,\delta K
    \;\sim\;
    \mathcal N\!\left(0,\;K+\delta K\right).
    \label{eq:xi_conditional}
\end{equation}
The reduced equilibrium operator then takes the nested stochastic form
\begin{align}
    \bar\rho_Q(\beta)
    &=
    Z_X(\beta)\,
    \mathbb E_{\delta K}\!\left[
        \mathbb E_{\xi|\delta K}\!\left[ U_\xi(\beta) \right]
    \right],
    \label{eq:nested_quenched}
\end{align}
where $\bar H(\tau;\xi)=H_Q-\xi(\tau)f$ exactly as in
\eqref{eq:bar_rho_as_quenched_average}. The inner average is the
familiar quenched construction of Sec.~\ref{sec:quenched}; the outer
average over the kernel fluctuation $\delta K$ captures the
non-Gaussian bath statistics. A Gaussian bath corresponds to the
special case $\delta K=0$ identically, and the hierarchy terminates at
a single level.

Crucially, in the commuting sector $[H_Q,f]=0$, the inner
Hubbard--Stratonovich average can be performed analytically---exactly as
in Sec.~\ref{sec:influence_to_hmf}---for every realisation of
$\delta K$. The result is
$e^{-\beta H_Q}\,\exp\!\bigl[\frac{1}{2}(C+\delta C)\,f^2\bigr]$,
where $\delta C:=\!\iint\!\delta K(\tau,\tau')\,d\tau\,d\tau'$ is a
scalar Gaussian variable with mean zero and variance
$8\,\kappa_4(\beta)$. The remaining outer average over $\delta C$ is
also Gaussian and evaluates in closed form:
\begin{align}
    \mathbb E_{\delta C}\!\left[e^{\frac{1}{2}\delta C\,f^2}\right]
    &=
    \exp\!\left[\tfrac{1}{2}\cdot\tfrac{1}{4}\cdot 8\kappa_4\,f^4\right]
    \nonumber\\
    &=
    \exp\!\left[\kappa_4(\beta)\,f^4\right],
    \label{eq:nested_commuting_check}
\end{align}
reproducing~\eqref{eq:HMF_non_gaussian} exactly. Both levels of the
hierarchy are evaluated analytically; no stochastic sampling is
required. The same pattern extends to higher cumulants: each
$\kappa_{2n}$ correction to the mean-force Hamiltonian can be obtained
by introducing one additional layer of Gaussian kernel fluctuation and
performing the resulting average in closed form.

For non-commuting couplings, the inner average can no longer be done
analytically, but the nested representation~\eqref{eq:nested_quenched}
still provides a constructive framework: sample $\delta K$, then sample
$\xi$ conditioned on $\delta K$, and average the quenched propagators
numerically. The construction requires $K+\delta K$ to be positive
semi-definite for each realisation, which holds perturbatively whenever
$\delta K\ll K$, i.e.\ when the fourth cumulant is small relative to
$K^{\otimes 2}$. For sub-Gaussian baths---including bounded spin baths,
where $K_4$ is typically negative---the HS identity remains valid as
an analytic continuation (the auxiliary field $\delta K$ becomes complex
valued), but the direct stochastic sampling interpretation is lost.

This perspective makes precise the two categories of obstruction to a
compact Hamiltonian of mean force. Gaussian bath statistics reduce the
cumulant series to a single kernel, while commutativity of $H_Q$ and
$f$ collapses the time ordering. The commuting Gaussian case treated in
Sec.~\ref{sec:influence_to_hmf} is the unique intersection where both
simplifications hold simultaneously. When either is relaxed, the nested
Hubbard--Stratonovich construction provides a systematic route back to
the local operator picture.
